\documentclass[10pt,a4paper]{amsart}
\usepackage[margin=19mm]{geometry}
\usepackage{amsmath,amssymb,mathtools}
\usepackage[scr=rsfs,cal=euler]{mathalfa}
\usepackage{xfrac}
\usepackage[unicode,hidelinks,bookmarks=false]{hyperref}
\newcommand{\ie}{i.e.\ }
\newcommand{\cf}{cf.\ }
\newcommand{\resp}{respectively\ }
\newcommand{\Subsection}{\subsection}
\providecommand{\nobreakdash}{\nobreak\hbox{-}}
\setlength{\emergencystretch}{2em}
\multlinegap0pt
\usepackage{amssymb}
\usepackage{xcolor}
\usepackage{tikz}
\newtheorem{theorem}{Theorem}
\newcommand{\LKb}{\mathbf{LK_b}}
\title{Exponential Gaps Between Intuitionistic Linear Extended Frege Systems}
\author{Amirhossein Akbar Tabatabai}
\date{Source: arXiv:2609.00422v1 (CC BY 4.0)}
\begin{document}
\maketitle

\noindent\emph{Source and licence.} Exponential Gaps Between Intuitionistic Linear Extended Frege Systems. Version arXiv:2609.00422v1 (CC BY 4.0), distributed by arXiv under the Creative Commons Attribution 4.0 International licence, http://creativecommons.org/licenses/by/4.0/, as recorded in the arXiv licence metadata for this exact version and shown on the abstract page https://arxiv.org/abs/2609.00422v1. Full licence notice: \texttt{licenses/arxiv-2609.00422-CC-BY-4.0.txt}.\par\smallskip

\noindent\textit{Editorial scope.} Counted unit: the article's theorem UnitPropagation (source0.tex lines 1111--1113). The article's complete original proof of it is delivered with it (source0.tex lines 1114--1237, one \texttt{proof} environment of 124 lines). No mathematics is written, repaired or re-derived: the statement and the proof are the article's own lines, in the article's own notation, and no retained line is substituted or re-flowed.  No further source-local context is needed: every cross-reference of every retained line resolves inside the two delivered spans.  Nothing was omitted from any retained span, so the package carries no omission record.  No retained line contains a citation, so the wrapper carries no bibliography and no bibliography entry.  No retained line contains a figure environment, an image-inclusion command or drawing code, so no image is delivered and the package declares no asset dependency.  Editorial formatting only, in addition: an AMS \texttt{amsart} preamble that restates the article's own declarations and macros used by the retained lines, namely the environments \texttt{theorem} on separate counters, together with the standard packages those lines need.  The retained environments print in this document as Theorem 1 in order of appearance, whereas the article prints the same items with the numbering of its own full document.  The complete native source of the article as distributed in the arXiv e-print for this exact version is bundled unchanged as \texttt{sources/209025/source0.tex}.
\begin{theorem}(Unit Propagation)\label{UnitPropagation}
For any set $\mathcal{S}$ of Horn sequents over $\mathcal{L}_b$ and any two atoms $p$ and $q$, if $\mathcal{S} \vdash_{\LKb} \, \Rightarrow p \vee q$, then there is $\tau$ such that $l(\tau) \leq |\mathcal{S}|^{O(1)}$ and either $\mathcal{S} \vdash_{\mathbf{FL_e}}^{\tau} \, \Rightarrow p$ or $\mathcal{S} \vdash_{\mathbf{FL_e}}^{\tau} \, \Rightarrow q$. 
\end{theorem}

\begin{proof}
Call any sequent in the form $(\, \Rightarrow u)$, where $u$ is an atom, a \emph{unit}. Now, define two sequences $\{\mathcal{S}_i\}_{i=0}^{\infty}$ and $\{\mathcal{U}_i\}_{i=0}^{\infty}$ of sequents in the following inductive way. Set $\mathcal{S}_0=\mathcal{S}$ and $\mathcal{U}_0=\varnothing$. Then, define $\mathcal{S}_{i+1}$ and $\mathcal{U}_{i+1}$ by applying the following process.
If there is no unit in $\mathcal{S}_i-\mathcal{U}_i$, set $\mathcal{S}_{i+1}=\mathcal{S}_{i}$ and $\mathcal{U}_{i+1}=\mathcal{U}_i$. Otherwise, pick a unit $(\, \Rightarrow u)$ arbitrarily in $\mathcal{S}_i-\mathcal{U}_i$. Now, for any sequent $S \in \mathcal{S}_i$ other than $(\, \Rightarrow u)$, apply the following simplifications to modify $\mathcal{S}_i$:
\begin{itemize}
\item[$(i)$] 
If $u$ does not appear in $S$, do not change $S$.
\item[$(ii)$] 
If $S$ is in the form $C_1, \cdots, C_m \Rightarrow u$, erase $S$ from $\mathcal{S}_i$.
\item[$(iii)$] 
If $S$ is in the form $C_1, \cdots, C_m \Rightarrow r$, where $r \neq u$, then substitute all occurrences of $u$ in the $C_i$'s by $1$.
\end{itemize}
When this process has been applied to all sequents of $\mathcal{S}_i$ except for $(\, \Rightarrow u)$, eliminate any formula that is just a conjunction of $1$'s from the antecedents of the resulting sequents, and set $\mathcal{S}_{i+1}$ to be the result. For $\mathcal{U}_{i+1}$, simply set it to $\mathcal{U}_i \cup \{\, \Rightarrow u\}$.

Here are some properties of these two sequences. First,
$\mathcal{U}_i \subseteq \mathcal{S}_i$ for any $i \in \mathbb{N}$ and each $\mathcal{S}_i$ is a set of Horn sequents satisfying $|\mathcal{S}_i|\subseteq |\mathcal{S}|$. To prove the first part, we use induction on $i \in \mathbb{N}$. For $i=0$, the claim is obvious. For the induction step, if there is no unit in $\mathcal{S}_i-\mathcal{U}_i$, then $\mathcal{S}_{i+1}=\mathcal{S}_{i}$ and $\mathcal{U}_{i+1}=\mathcal{U}_i$, and the claim follows from the induction hypothesis for $i$. Otherwise, let $(\, \Rightarrow u)$ be the chosen unit in $\mathcal{S}_i-\mathcal{U}_i$. By definition, $(\, \Rightarrow u)$ belongs to $\mathcal{S}_i$ and remains unchanged throughout the modifications. Hence, $(\, \Rightarrow u) \in \mathcal{S}_{i+1}$. For any other unit $(\, \Rightarrow v) \in \mathcal{U}_{i+1}$, it must belong to $\mathcal{U}_i$ and hence, by the induction hypothesis, $(\, \Rightarrow v) \in \mathcal{S}_i$. As $v \neq u$, the sequent $(\, \Rightarrow v)$ remains intact throughout the modifications, which only affect occurrences of $u$. Therefore, $(\, \Rightarrow v) \in \mathcal{S}_{i+1}$. For the second part, being Horn is clear from the modifications. For the bound, it is enough to note that $|\mathcal{S}_{i+1}| \leq |\mathcal{S}_i|$, for any $i\in \mathbb{N}$. 

Second, by an easy induction, one can show that, for any $i\in\mathbb{N}$, if an atom occurs in $\mathcal{U}_i$, then it appears only in a unit in $\mathcal{S}_i$. For $i=0$, there is nothing to prove. For the induction step, if there is no unit in $\mathcal{S}_i-\mathcal{U}_i$, then $\mathcal{S}_{i+1}=\mathcal{S}_i$ and $\mathcal{U}_{i+1}=\mathcal{U}_i$, and the claim follows from the induction hypothesis for $i$. Otherwise, let $(\,\Rightarrow u)$ be the chosen unit in $\mathcal{S}_i-\mathcal{U}_i$, and let $(\,\Rightarrow v)$ be a unit in $\mathcal{U}_{i+1}$. There are two cases. If $v=u$, then, since all occurrences of $u$ in $\mathcal{S}_i$ except for $(\,\Rightarrow u)$ are eliminated when defining $\mathcal{S}_{i+1}$, the only occurrence of $u$ in $\mathcal{S}_{i+1}$ is in the unit $(\,\Rightarrow u)$. If $v\neq u$, then $(\,\Rightarrow v)\in\mathcal{U}_i$, and by the induction hypothesis, $v$ occurs only in units of $\mathcal{S}_i$. Since the modifications only remove sequents or replace occurrences of $u$, the claim follows.

Third, we claim 
\[
\LKb \vdash
\bigwedge_{S\in\mathcal{S}_{i+1}} I(S)
\Leftrightarrow
\bigwedge_{S\in\mathcal{S}_{i}} I(S),
\]
for any $i\in\mathbb{N}$. This follows directly from the following
$\LKb$-equivalences:
\begin{itemize}
    \item
    $u\wedge(A\to u)\Leftrightarrow u$;
    
    \item
    $u\wedge((A\wedge u)*B\to r)
    \Leftrightarrow
    u\wedge((A\wedge 1)*B\to r)$;
    
    \item
    $(A*B\to r)\Leftrightarrow(B\to r)$,
    where $A$ is a conjunction of $1$'s.
\end{itemize}
Notice that we are working in the system $\LKb$, where all structural rules
are available. Hence, these are simply the usual classical equivalences
expressed in the substructural language.

Fourth, for any $i\in\mathbb{N}$, there is an $\mathbf{FL_e}$-proof with
$|\mathcal{S}|^{O(1)}$ many lines for
\[
\Rightarrow\bigwedge_{S\in\mathcal{S}_i} I(S)
\vdash_{\mathbf{FL_e}}
\Rightarrow\bigwedge_{S\in\mathcal{S}_{i+1}} I(S).
\]
There are two cases to consider. If there is no unit in
$\mathcal{S}_i-\mathcal{U}_i$, then
$\mathcal{S}_{i+1}=\mathcal{S}_i$, and hence there is nothing to prove.
Otherwise, let $(\,\Rightarrow u)\in\mathcal{S}_i-\mathcal{U}_i$ be the
chosen unit and let $T$ be an arbitrary sequent in $\mathcal{S}_{i+1}$.
Assume that $T$ is obtained by modifying some $R\in\mathcal{S}_i$.
Since conjunctions of $1$'s are equivalent to $1$ and can be eliminated
from antecedents in $\mathbf{FL_e}$, without loss of generality, we may
assume that $R$ is the result of either $(i)$ or $(iii)$. Note that $(ii)$
does not apply, as it eliminates sequents.
For $(i)$, we have $T=R\in\mathcal{S}_i$, and hence
$
\Rightarrow\bigwedge_{S\in\mathcal{S}_i} I(S)
\vdash_{\mathbf{FL_e}}
\Rightarrow I(T).
$
For $(iii)$, let
\[
R=(\Gamma,A_1\wedge u,\ldots,A_m\wedge u\Rightarrow r),
\]
where $r\neq u$ and $\Gamma$ contains no occurrence of $u$. Then
\[
T=(\Gamma,A_1\wedge1,\ldots,A_m\wedge1\Rightarrow r).
\]
Now, observe that
\[
\{(\,\Rightarrow u),
(\Gamma,\{A_j\wedge u\}_{j=1}^m\Rightarrow r)\}
\vdash_{\mathbf{FL_e}}
(\Gamma,\{A_j\wedge1\}_{j=1}^m\Rightarrow r).
\]
Therefore, since $(\,\Rightarrow u)\in\mathcal{S}_i$, we obtain
$
\Rightarrow\bigwedge_{S\in\mathcal{S}_i} I(S)
\vdash_{\mathbf{FL_e}}
\Rightarrow I(T).
$
For the number of lines in the proof, first note that for each
$T\in\mathcal{S}_{i+1}$ in case $(i)$, the number of lines is clearly
$O(|\mathcal{S}_i|)\leq O(|\mathcal{S}|)$, by the first property. In case
$(iii)$, the number of lines is
$|\mathcal{S}_i|^{O(1)}\leq|\mathcal{S}|^{O(1)}$, again by the first
property. Since the number of sequents $T$ is
$O(|\mathcal{S}_{i+1}|)\leq O(|\mathcal{S}|)$, we obtain an
$\mathbf{FL_e}$-proof with $|\mathcal{S}|^{O(1)}$ many lines witnessing
\[
\Rightarrow\bigwedge_{S\in\mathcal{S}_i} I(S)
\vdash_{\mathbf{FL_e}}
\Rightarrow\bigwedge_{S\in\mathcal{S}_{i+1}} I(S).
\]
Now, we use these four properties to prove the claim. First, by the first
property, we have $\mathcal{U}_i\subseteq\mathcal{S}_i$ and
$|\mathcal{S}_i|\leq|\mathcal{S}|$ for any $i\in\mathbb{N}$. We claim that
at some point $N\in\mathbb{N}$, there is no unit in
$\mathcal{S}_N-\mathcal{U}_N$. Otherwise, the sequence
$\{\mathcal{U}_i\}_{i\in\mathbb{N}}$ would be a strictly increasing
sequence of sets whose sizes are bounded by $|\mathcal{S}|$, which is
impossible. Now, let $N$ be a stage at which the two sequences
$\{\mathcal{S}_i\}_{i\in\mathbb{N}}$ and
$\{\mathcal{U}_i\}_{i\in\mathbb{N}}$ stabilize. Notice that $N$ is bounded
by the number of atoms occurring in $\mathcal{S}$, and hence
$N\leq|\mathcal{S}|$.

We claim that either $(\, \Rightarrow p)$ or $(\, \Rightarrow q)$ appears in $\mathcal{S}_N$. Let us assume otherwise. As $\mathcal{S} \vdash_{\LKb} \, \Rightarrow p \vee q$, we have $\LKb \vdash \bigwedge_{S \in \mathcal{S}} I(S) \Rightarrow p \vee q$. Now, by the third property, we reach $\LKb \vdash \bigwedge_{S \in \mathcal{S}_N} I(S) \Rightarrow p \vee q$, which is equivalent to $\mathcal{S}_N \vdash_{\LKb} \, \Rightarrow p \vee q$. Therefore, any Boolean valuation satisfying $\mathcal{S}_N$ must validate $(\, \Rightarrow p \vee q)$.
Define the valuation $V$ by setting $V(u)=1$ if and only if $(\, \Rightarrow u) \in \mathcal{U}_N$. We claim that $V(S)=1$ for any $S \in \mathcal{S}_N$. If $S$ is a unit, then $S=(\, \Rightarrow u)$. As there is no unit in $\mathcal{S}_N-\mathcal{U}_N$, we must have $(\, \Rightarrow u) \in \mathcal{U}_N$. Therefore, $V(u)=1$, which implies $V(S)=1$. If $S$ is not a unit, then $S=(C_1, \cdots, C_m \Rightarrow r)$, where $m$ is non-zero. Therefore, at least one atom $s$ appears in the antecedent of $S$. By the second property, since $s$ appears in a non-unit sequent $S \in \mathcal{S}_N$, we must have $(\, \Rightarrow s) \notin \mathcal{U}_N$. Hence, $V(s)=0$, which implies $V(S)=1$. Therefore, $V$ validates all sequents in $\mathcal{S}_N$.

Now, by the first property, $\mathcal{U}_i \subseteq \mathcal{S}_i$ for any $i \in \mathbb{N}$. Hence, by the assumption that $(\, \Rightarrow p)$ and $(\, \Rightarrow q)$ are not in $\mathcal{S}_N$, they are also outside of $\mathcal{U}_N$. Thus, $V(p)=V(q)=0$. Therefore, $V$ does not validate $(\, \Rightarrow p \vee q)$, which is a contradiction with $\mathcal{S}_N \vdash_{\LKb} \, \Rightarrow p \vee q$. Therefore, either $(\, \Rightarrow p)$ or $(\, \Rightarrow q)$ appears in $\mathcal{S}_N$. Hence, either $\Rightarrow \bigwedge_{S \in \mathcal{S}_N} I(S) \vdash_{\mathbf{FL_e}} \, \Rightarrow p$ or $\Rightarrow \bigwedge_{S \in \mathcal{S}_N} I(S) \vdash_{\mathbf{FL_e}} \, \Rightarrow q$. Note that the proofs have $O(|\mathcal{S}_N|)\leq O(|\mathcal{S}|)$ many lines.
By the fourth property, for any $i \in \mathbb{N}$, we have 
\[
\, \Rightarrow \bigwedge_{S \in \mathcal{S}_i} I(S) \vdash_{\mathbf{FL_e}} \, \Rightarrow \bigwedge_{S \in \mathcal{S}_{i+1}} I(S)
\]
with a proof containing $|\mathcal{S}|^{O(1)}$ many lines. As $N \leq |\mathcal{S}|$, we obtain a proof of either $\, \Rightarrow \bigwedge_{S \in \mathcal{S}} I(S) \vdash_{\mathbf{FL_e}} \, \Rightarrow p$ or $\, \Rightarrow \bigwedge_{S \in \mathcal{S}} I(S) \vdash_{\mathbf{FL_e}} \, \Rightarrow q$ with $|\mathcal{S}|^{O(1)}$ many lines. Therefore, there is an $\mathbf{FL_e}$-proof $\tau$ of $\mathcal{S} \vdash_{\mathbf{FL_e}} \, \Rightarrow p$ or $\mathcal{S} \vdash_{\mathbf{FL_e}} \, \Rightarrow q$ such that $l(\tau) \leq |\mathcal{S}|^{O(1)}$.
\end{proof}

\end{document}
